% ============================================================
%  Abstract & Keywords
% ============================================================
\begin{abstract}
  In large-scale rotating machinery such as turbines, generators,
  and compressors, shaft alignment accuracy is critical to
  operational balance and equipment lifespan.
  %
  Traditional steel-wire alignment methods suffer from
  wire sag, vibration sensitivity, and limited measurement range.
  %
  This paper proposes DCPAM
  (Dual-Camera-based Point-to-Axis Measurement),
  a non-contact measurement model that replaces the steel wire
  with a laser reference line and employs a dual-camera system
  to reconstruct the three-dimensional laser axis.
  %
  The model computes the perpendicular distance from any target
  point to the reconstructed laser axis through a pipeline of
  pixel extraction, camera-to-world coordinate transformation,
  mirror-reflection virtual-image mapping, and cross-product
  distance calculation.
  %
  Two parameter-tuning strategies are presented:
  DCPAM-CV, an experience-driven classical computer-vision
  workflow, and DCPAM-CNN, a convolutional-neural-network-based
  approach.
  %
  A closed-form error propagation analysis and Monte Carlo
  validation quantify the measurement uncertainty under
  micrometer-level perturbations.
  % TODO: 补充实验结果摘要——最终精度数据
\end{abstract}

\begin{IEEEkeywords}
  Shaft alignment, laser measurement, dual-camera system,
  point-to-axis distance, error propagation, computer vision.
\end{IEEEkeywords}
